\documentclass[11pt]{article}
\usepackage{worksheet_style}

% ---- Toggle: uncomment for filled version ----
%\worksheetfilledtrue
\begin{document}
\worksheetheader{13}{Double Integrals over Rectangular Regions}

\begin{background}
We now enter the second major part of the course: integration in multiple dimensions. Just as the single integral $\int_a^b f(x)\,dx$ computes area under a curve, the double integral $\iint_R f(x,y)\,dA$ computes volume under a surface. Today we start with rectangular regions and Fubini's Theorem, which lets us evaluate double integrals as iterated single integrals. This is the gateway to computing volumes, masses, centers of mass, and much more over two- and three-dimensional regions.
\end{background}

%=============================================================
\wsection{Double Integrals as Volume via Riemann Sums}

\begin{defbox}{Double Integral}
The \textbf{double integral} of $f(x,y)$ over a rectangle $R = [a,b]\times[c,d]$:
\[\iint_R f(x,y)\,dA = \wbml{\lim_{m,n\to\infty}\sum_{i=1}^{m}\sum_{j=1}^{n} f(x_{ij}^*, y_{ij}^*)\,\Delta A}\]

-- Partition $R$ into $m \times n$ subrectangles, each with area $\Delta A = \Delta x \cdot \Delta y$

-- Pick a sample point $(x_{ij}^*, y_{ij}^*)$ in each subrectangle

-- If $f \geq 0$: the double integral $=$ \wbi{volume under $z = f(x,y)$ above $R$}{7cm}
\end{defbox}

\begin{center}
\tdplotsetmaincoords{70}{110}
\begin{tikzpicture}[tdplot_main_coords,line cap=round,>=stealth,
    declare function={f(\x,\y)=3+0.5*sin(30*\x)*cos(30*\y);}]
 \draw[->] (0,0,0) coordinate (O) -- (5,0,0) node[pos=1.05]{$x$};
 \draw[->] (O) -- (0,5,0) node[pos=1.05]{$y$};
 \draw[->] (O) -- (0,0,5) node[pos=1.05]{$z$};
 \draw (1,2,0) -- (2,2,0) -- (2,3,0) -- (1,3,0) -- cycle;
 \foreach \X in {1,2} {\foreach \Y in {2,3}
   {\draw[dashed] (\X,\Y,0) -- (\X,\Y,{f(\X,\Y)});}}
 \foreach \X in {1,2} {
   \draw[dashed] (\X,2,0) -- (\X,0,0);}
 \foreach \Y in {2,3} {
   \draw[dashed] (1,\Y,0) -- (0,\Y,0);}
 \foreach \X in {0,...,4}
  {\draw plot[variable=\x,domain=0:5,smooth] (\X,\x,{f(\X,\x)});}
 \foreach \Y in {0,...,5}
  {\draw plot[variable=\x,domain=0:4,smooth] (\x,\Y,{f(\x,\Y)});}
 \draw[blue,thick] plot[variable=\x,domain=1:2] (\x,2,{f(\x,2)})
    -- plot[variable=\y,domain=2:3] (2,\y,{f(2,\y)})
    -- plot[variable=\x,domain=2:1] (\x,3,{f(\x,3)})
    -- plot[variable=\y,domain=3:2] (1,\y,{f(1,\y)})
    -- cycle;
\end{tikzpicture}\\
{\small The double integral extends the idea of area under a curve to volume under a surface. Each sub-rectangle in the $xy$-plane gets a little prism of height $f(x,y)$; adding them up in the limit gives $\iint_R f\,dA$.}
\end{center}

\begin{example}[Riemann Sum Approximation]
Approximate $\ds\iint_R (x^2 + y^2)\,dA$ over $R = [0,2]\times[0,1]$ using a $2\times 2$ grid with upper-right sample points.

1. \textbf{Set up grid:}

$\Delta x = $ \wbm{$1$}, \quad $\Delta y = $ \wbm{$0.5$}, \quad $\Delta A = $ \wbm{$0.5$}

Upper-right sample points: $(1, 0.5)$, $(2, 0.5)$, $(1, 1)$, $(2, 1)$

2. \textbf{Evaluate $f$ at each sample point:}

\wbbox{
$f(1, 0.5) = 1 + 0.25 = 1.25$

$f(2, 0.5) = 4 + 0.25 = 4.25$

$f(1, 1) = 1 + 1 = 2$

$f(2, 1) = 4 + 1 = 5$
}{3cm}

3. \textbf{Compute sum:}

$\ds\sum f \cdot \Delta A = (1.25 + 4.25 + 2 + 5)(0.5) = 12.5 \times 0.5 = $ \wbm{$6.25$}

\medskip
\textit{Exact answer (computed below): $10/3 \approx 3.33$. The Riemann sum overestimates because $f$ is increasing and we used upper-right points.}
\end{example}

%=============================================================
\wsection{Fubini's Theorem and Iterated Integrals}

\begin{thmbox}{Fubini's Theorem}
If $f$ is continuous on $R = [a,b]\times[c,d]$, then:
\[\iint_R f(x,y)\,dA = \int_a^b\!\int_c^d f(x,y)\,dy\,dx = \int_c^d\!\int_a^b f(x,y)\,dx\,dy\]

-- The order of integration does not matter for continuous functions on rectangles

-- Evaluate from the \textbf{inside out}: inner integral first, then outer
\end{thmbox}

\begin{example}[Iterated Integral of $xy$]
Evaluate $\ds\iint_R xy\,dA$ where $R = [0,3]\times[0,2]$.

1. \textbf{Set up (integrate $y$ first):}

$\ds\int_0^3\!\int_0^2 xy\,dy\,dx$

2. \textbf{Inner integral} (treat $x$ as constant):

$\ds\int_0^2 xy\,dy = \left[\frac{xy^2}{2}\right]_0^2 = $ \wbm{$2x$}

3. \textbf{Outer integral:}

$\ds\int_0^3 2x\,dx = [x^2]_0^3 = $ \wbm{$9$}
\end{example}

\begin{example}[Symmetry Shortcut]
Evaluate $\ds\iint_R \frac{y}{x^2+y^2}\,dA$ where $R = [1,2]\times[-1,1]$.

\wbbox{
\textbf{Key insight:} $f(x,y) = \frac{y}{x^2+y^2}$ is an \textbf{odd function of $y$}: $f(x,-y) = -f(x,y)$.

The region $R$ is symmetric about the $x$-axis ($y$ ranges from $-1$ to $1$).

By symmetry: $\ds\iint_R \frac{y}{x^2+y^2}\,dA = \boxed{0}$

No computation needed!
}{3cm}
\end{example}

\begin{formulabox}{Symmetry Principle}
If $f$ is \textbf{odd} in $y$ (i.e., $f(x,-y) = -f(x,y)$) and $R$ is symmetric about the $x$-axis:
\[\iint_R f(x,y)\,dA = 0\]
Same idea for odd in $x$ on a region symmetric about the $y$-axis.
\end{formulabox}

%=============================================================
\wsection{Strategic Order of Integration}

-- Sometimes one order is much easier than the other

-- Look for which order avoids difficult antiderivatives

\begin{example}[Computational Iterated Integral]
Evaluate $\ds\int_0^1\!\int_0^2 (x^3\sqrt{y} + \cos x)\,dy\,dx$.

1. \textbf{Inner integral} (integrate with respect to $y$):

\wbbox{
$\ds\int_0^2 (x^3 y^{1/2} + \cos x)\,dy = \left[\frac{2x^3}{3}y^{3/2} + y\cos x\right]_0^2 = \frac{2x^3}{3}\cdot 2\sqrt{2} + 2\cos x = \frac{4\sqrt{2}}{3}x^3 + 2\cos x$
}{2.5cm}

2. \textbf{Outer integral:}

\wbbox{
$\ds\int_0^1 \left(\frac{4\sqrt{2}}{3}x^3 + 2\cos x\right)dx = \frac{4\sqrt{2}}{3}\cdot\frac{x^4}{4}\Big|_0^1 + 2\sin x\Big|_0^1 = \frac{\sqrt{2}}{3} + 2\sin 1$
}{2cm}

3. \textbf{Result:} $\boxed{\dfrac{\sqrt{2}}{3} + 2\sin 1}$
\end{example}

%=============================================================

\begin{practice}

\prob Evaluate $\ds\int_0^1\!\int_0^2 (x + 2y)\,dx\,dy$.

\wbbox{
$\int_0^1\left[\frac{x^2}{2} + 2xy\right]_0^2\,dy = \int_0^1(2 + 4y)\,dy = [2y + 2y^2]_0^1 = 4$.
}{2.5cm}

\wans{\wbm{$4$}}

\vspace{2.5cm}

\prob Evaluate $\ds\iint_R e^{x+y}\,dA$ where $R = [0,1]\times[0,1]$.

\wbbox{
$\int_0^1\!\int_0^1 e^x e^y\,dy\,dx = \int_0^1 e^x[e^y]_0^1\,dx = (e-1)\int_0^1 e^x\,dx = (e-1)[e^x]_0^1 = (e-1)^2$.
}{2.5cm}

\wans{\wbm{$(e-1)^2$}}

\vspace{2.5cm}

\prob Verify Fubini's Theorem by evaluating $\ds\iint_R 2xy\,dA$ over $R = [1,3]\times[0,2]$ in both orders.

\wbbox{
Order 1: $\int_1^3\!\int_0^2 2xy\,dy\,dx = \int_1^3[xy^2]_0^2\,dx = \int_1^3 4x\,dx = [2x^2]_1^3 = 18-2 = 16$.

Order 2: $\int_0^2\!\int_1^3 2xy\,dx\,dy = \int_0^2[x^2 y]_1^3\,dy = \int_0^2 8y\,dy = [4y^2]_0^2 = 16$. \;\checkmark
}{3.5cm}

\wans{\wbm{$16$} (both orders agree)}

\vspace{3cm}

\prob Without computing, explain why $\ds\iint_{[-1,1]\times[-1,1]} x^3 y\,dA = 0$.

\wbbox{
$f(x,y) = x^3 y$ is odd in $x$: $f(-x,y) = -x^3 y = -f(x,y)$.

The region $[-1,1]\times[-1,1]$ is symmetric about the $y$-axis.

By symmetry, the contributions from $x > 0$ and $x < 0$ cancel exactly, giving $0$.
}{2.5cm}

\wans{\wbm{$0$} (by symmetry --- $f$ is odd in $x$ on a symmetric region)}

\vspace{2.5cm}

\prob Evaluate $\ds\iint_R (x^2 y - y^2)\,dA$ where $R = [0,2]\times[-1,1]$.

\wbbox{
Split: $\iint_R x^2 y\,dA - \iint_R y^2\,dA$.

First integral: $x^2 y$ is odd in $y$ on $[-1,1]$, so $\iint_R x^2 y\,dA = 0$.

Second integral: $\int_0^2\!\int_{-1}^1 y^2\,dy\,dx = \int_0^2\left[\frac{y^3}{3}\right]_{-1}^1 dx = \int_0^2 \frac{2}{3}\,dx = \frac{4}{3}$.

Total: $0 - \frac{4}{3} = -\frac{4}{3}$.
}{3.5cm}

\wans{\wbm{$-4/3$}}

\end{practice}

\end{document}
